% Chapter 6: Sard's Theorem
% Lee, Introduction to Smooth Manifolds (2nd ed., GTM 218).
% Generated from the section extraction; nodes carry only Lean that is
% verified sorry-free. See ../../UPSTREAM_LEAN_AUDIT.md.


\section{Sets of Measure Zero}

\begin{definition}
\label{def:6-38-extra-1}
\lean{NullMeasurableSet.of_null}

\mathlibok
A set $A \subseteq \mathbb{R}^n$ has measure zero if, for every $\delta>0$, there is a countable collection of open rectangles covering $A$ whose volumes sum to less than $\delta$.
\end{definition}

\begin{lemma}[Equivalent Null-Cover Formulations]
\label{exer:6-1}
\lean{volume_eq_zero_iff_forall_pos_exists_open_cube_cover}


\dcref{ch6:6.1}
The definition of measure zero is unchanged if open rectangles are replaced
by open balls or by open cubes. Every open rectangle
$R\subseteq\mathbb R^n$ has a finite open-cube cover of total volume at most
$2^n$ times the volume of $R$.
\end{lemma}

\begin{lemma}
\label{lem:6-2}
\lean{volume_eq_zero_of_isCompact_of_forall_first_coordinate_slice_volume_eq_zero}
\uses{exer:6-1}
\leanok



\dcref{ch6:6.2}
Let $A \subseteq \mathbb{R}^n$ be compact. If
$A \cap (\{c\}\times\mathbb{R}^{n-1})$ has $(n-1)$-dimensional measure zero for every $c\in\mathbb{R}$, then $A$ has $n$-dimensional measure zero.
\end{lemma}

\begin{proof}
\leanok
\sketch Choose $[a,b]$ with $A\subseteq [a,b]\times\mathbb{R}^{n-1}$ and set
\[
A_c=\{x\in\mathbb{R}^{n-1}:(c,x)\in A\}.
\]
For $\delta>0$, cover each compact null set $A_c$ by finitely many open cubes with total volume less than $\delta$, and denote their union by $U_c$. Compactness of $A$ gives an interval $J_c$ about $c$ such that
\[
A\cap(J_c\times\mathbb{R}^{n-1})\subseteq J_c\times U_c;
\]
otherwise a convergent sequence in $A$ would yield a point of $A_c\setminus U_c$. Choose finitely many $J_{c_i}$ covering $[a,b]$ and shrink overlaps so that their total length is at most $2|b-a|$. The corresponding rectangles cover $A$ and have total volume less than $2\delta|b-a|$. Letting $\delta$ tend to zero proves the claim.
\end{proof}

\begin{proposition}
\label{prop:6-3}
\lean{volume_halfSpace_graph_eq_zero_of_measurableSet_of_continuous}
\uses{lem:6-2}
\leanok



\dcref{ch6:6.3}
If $A$ is an open or closed subset of $\mathbb{R}^{n-1}$ or $\mathbb{H}^{n-1}$ and $f:A\to\mathbb{R}$ is continuous, then the graph of $f$ has measure zero in $\mathbb{R}^n$.
\end{proposition}

\begin{proof}
\leanok
\sketch For compact $A$, argue by induction on $n$. The case $n=1$ is a set with at most one point. For each $c\in\mathbb{R}$, the slice of the graph in $\{c\}\times\mathbb{R}^{n-1}$ is the graph of a continuous function of $n-2$ variables, hence is null by induction. Lemma~\ref{lem:6-2} makes the entire graph null. In general, Proposition A.60 expresses $A$ as a countable union of compact sets, so its graph is a countable union of null sets.
\end{proof}

\begin{corollary}
\label{cor:6-4}
\lean{volume_affineSubspace_eq_zero_of_ne_top}
\uses{prop:6-3}
\leanok



\dcref{ch6:6.4}
Every proper affine subspace of $\mathbb{R}^n$ has measure zero in $\mathbb{R}^n$.
\end{corollary}

\begin{proof}
\leanok
\sketch If $\dim S=n-1$, choose a coordinate axis not parallel to $S$. Then $S$ is the graph of an affine function in the remaining coordinates and is null by Proposition~\ref{prop:6-3}. If $\dim S<n-1$, embed $S$ in an affine hyperplane and apply Proposition C.18(b).
\end{proof}

\begin{definition}
\label{def:6-38-extra-2}
\lean{has_measure_zero_in_manifold}
\leanok


Let $M$ be a smooth $n$-manifold with or without boundary. A subset $A\subseteq M$ has measure zero in $M$ if $\varphi(A\cap U)$ has $n$-dimensional measure zero in $\mathbb{R}^n$ for every smooth chart $(U,\varphi)$ of $M$.
\end{definition}

\begin{lemma}
\label{lem:6-6}
\lean{has_measure_zero_in_manifold_of_chart_cover}
\leanok


\dcref{ch6:6.6}
Let $M$ be a smooth $n$-manifold with or without boundary and $A\subseteq M$. Suppose a collection of smooth charts $\{(U_\alpha,\varphi_\alpha)\}$ covers $A$ and each $\varphi_\alpha(A\cap U_\alpha)$ has measure zero in $\mathbb{R}^n$. Then $A$ has measure zero in $M$.
\end{lemma}

\begin{proof}
\leanok
\sketch Fix a smooth chart $(V,\psi)$. A countable subcollection of the $U_\alpha$ covers $A\cap V$. For every chart in this subcollection,
\[
\psi(A\cap V\cap U_\alpha)
=(\psi\circ\varphi_\alpha^{-1})
\bigl(\varphi_\alpha(A\cap V\cap U_\alpha)\bigr).
\]
The inner set is null by the hypothesis, and Proposition 6.5 shows that its image under the smooth transition map is null. Their countable union is $\psi(A\cap V)$, so this set is null.
\end{proof}

\begin{proposition}[Countable Stability of Null Sets]
\label{exer:6-7}
\lean{has_measure_zero_in_manifold_iUnion}
\leanok


\dcref{ch6:6.7}
A countable union of measure-zero subsets of a smooth manifold, with or
without boundary, has measure zero.
\end{proposition}

\begin{proposition}
\label{prop:6-8}
\lean{has_measure_zero_in_manifold}
\leanok


\dcref{ch6:6.8}
If $A$ has measure zero in a smooth manifold $M$ with or without boundary, then $M\setminus A$ is dense in $M$.
\end{proposition}

\begin{proof}
\leanok
\sketch Otherwise $A$ contains a nonempty open subset of $M$. In some smooth chart $(V,\psi)$, the null set $\psi(A\cap V)\subseteq\mathbb{R}^{\dim M}$ would then contain a nonempty open set, contradicting Corollary C.25.
\end{proof}

\begin{theorem}
\label{thm:6-9}
\lean{has_measure_zero_in_manifold}
\leanok


\dcref{ch6:6.9}
Let $M$ and $N$ be smooth $n$-manifolds with or without boundary. If $F:M\to N$ is smooth and $A\subseteq M$ has measure zero, then $F(A)$ has measure zero in $N$.
\end{theorem}

\begin{proof}
\leanok
\sketch Cover $M$ by countably many smooth charts $(U_i,\varphi_i)$. For each smooth chart $(V,\psi)$ of $N$,
\[
\psi(F(A)\cap V)
=\bigcup_i(\psi\circ F\circ\varphi_i^{-1})
\bigl(\varphi_i(A\cap U_i\cap F^{-1}(V))\bigr).
\]
Every set on the right is null by Proposition 6.5, and a countable union of null sets is null.
\end{proof}


\section{The Whitney Approximation Theorems}

\begin{definition}
\label{def:6-41-extra-1}
\lean{delta_close}
\leanok


For a positive continuous function $\delta:M\to\mathbb{R}$, maps $F,\widetilde F:M\to\mathbb{R}^k$ are $\delta$-close if
\[
|F(x)-\widetilde F(x)|<\delta(x)
\qquad\text{for every }x\in M.
\]
\end{definition}

\begin{corollary}
\label{cor:6-22}
\lean{Manifold.exists_positive_smooth_lt}
\leanok


\dcref{ch6:6.22}
If $M$ is a smooth manifold with or without boundary and $\delta:M\to\mathbb{R}$ is positive and continuous, then there is a smooth function $e:M\to\mathbb{R}$ such that
\[
0<e(x)<\delta(x)
\qquad\text{for every }x\in M.
\]
\end{corollary}

\begin{proof}
\leanok
\sketch Apply the Whitney approximation theorem to $\delta/2$ to obtain a smooth $e$ satisfying
\[
|e(x)-\delta(x)/2|<\delta(x)/2
\]
for every $x\in M$.
\end{proof}


\section{Tubular Neighborhoods}

\begin{definition}
\label{def:6-42-extra-1}
\lean{normal_bundle}
\leanok


Use the canonical identifications $T_x\mathbb{R}^n\cong\mathbb{R}^n$ and $T\mathbb{R}^n\cong\mathbb{R}^n\times\mathbb{R}^n$, with their Euclidean inner products. If $M\subseteq\mathbb{R}^n$ is an embedded $m$-submanifold, its normal space at $x\in M$ is the $(n-m)$-dimensional subspace
\[
N_xM=(T_xM)^\perp\subseteq T_x\mathbb{R}^n.
\]
Its normal bundle and projection are
\[
NM=\{(x,v)\in\mathbb{R}^n\times\mathbb{R}^n:x\in M,\ v\in N_xM\},
\qquad
\pi_{NM}=\pi|_{NM}:NM\to M.
\]
\end{definition}

\begin{definition}
\label{def:6-42-extra-3}
\lean{TopologicalRetraction}
\leanok


A retraction of a topological space $X$ onto $M\subseteq X$ is a continuous map $r:X\to M$ such that $r|_M=\id_M$.
\end{definition}


\section{Smooth Approximation of Maps Between Manifolds}

\begin{definition}
\label{def:6-43-extra-1}
\lean{ContMDiffMap.SmoothlyHomotopic}
\leanok


For smooth manifolds $N$ and $M$ with or without boundary, a homotopy $H:N\times I\to M$ is smooth if it extends smoothly to a neighborhood of $N\times I$ in $N\times\mathbb{R}$. Two maps are smoothly homotopic if they are joined by such a homotopy.
\end{definition}

\begin{lemma}
\label{lem:6-28}
\lean{ContMDiffMap.SmoothlyHomotopic.equivalence}
\leanok


\dcref{ch6:6.28}
Smooth homotopy is an equivalence relation on the smooth maps between any smooth manifolds $N$ and $M$ with or without boundary.
\end{lemma}

\begin{proof}
\leanok
\sketch Reflexivity and symmetry follow from constant homotopies and reversal of the interval. For transitivity, let $H_1$ join $F$ to $G$ and $H_2$ join $G$ to $K$. Choose a smooth $\varphi:[0,1]\to[0,2]$ such that $\varphi(0)=0$, $\varphi(1)=2$, $\varphi([0,1/2])\subseteq[0,1]$, $\varphi([1/2,1])\subseteq[1,2]$, and $\varphi=1$ near $1/2$. Define

\[
H(x,t)=\begin{cases}
H_1(x,\varphi(t)),&0\leq t\leq \tfrac12,\\
H_2(x,\varphi(t)-1),&\tfrac12\leq t\leq1.
\end{cases}
\]
Because both pieces equal $G$ near $t=1/2$, they glue smoothly; thus $H$ joins $F$ to $K$.
\end{proof}


\section{Transversality}

\begin{definition}
\label{def:6-44-extra-2}
\lean{IsSmoothFamily}
\leanok


Let $N$, $M$, and $S$ be smooth manifolds. A family of maps $\{F_s:N\to M\}_{s\in S}$ is smooth if the associated map
\[
F:N\times S\to M,
\qquad F(x,s)=F_s(x),
\]
is smooth.
\end{definition}

\begin{proposition}
\label{prop:6-34}
\lean{chartedSpace_locPathConnectedSpace}
\leanok


\dcref{ch6:6.34}
Let $\{F_s:N\to M\}_{s\in S}$ be a smooth family with $S$ connected. Then $F_{s_1}$ and $F_{s_2}$ are homotopic for every $s_1,s_2\in S$.
\end{proposition}

\begin{proof}
\leanok
\sketch A connected smooth manifold is path-connected. For a path $\gamma:[0,1]\to S$ from $s_1$ to $s_2$, the map
\[
H(x,t)=F(x,\gamma(t))
\]
is a homotopy from $F_{s_1}$ to $F_{s_2}$.
\end{proof}

\begin{definition}
\label{def:6-44-extra-3}
\lean{has_measure_zero_in_manifold}
\leanok


A subset $B$ of a smooth manifold $S$ contains almost every element of $S$ if $S\setminus B$ has measure zero in $S$.
\end{definition}
